\documentclass[11pt]{article}
\usepackage{amsmath,amssymb,amsthm}
\usepackage[margin=1in]{geometry}
\newtheorem{theorem}{Theorem}
\newtheorem{lemma}{Lemma}
\newtheorem{corollary}{Corollary}
\theoremstyle{definition}
\newtheorem{remark}{Remark}
\DeclareMathOperator{\Imag}{Im}
\newcommand{\Z}{\mathbb{Z}}
\newcommand{\Q}{\mathbb{Q}}
\newcommand{\F}{\mathbb{F}}

\title{Two-row $d=4$ law, classes $b\equiv 2,3 \pmod 4$:\\
the unique $2$-adic root and a single-candidate certificate}
\author{Clio}
\date{2026-06-09 (prove session)}

\begin{document}
\maketitle

\noindent\textbf{Status.} $(\star)$ \emph{not closed uniformly.} The bare-factor-$m$
``obstruction'' is shown to make $Q_b$ have a \emph{unique} root in $\Z_2$, so $Q_b$ has
\emph{at most one} rational root. This yields the cleanest known certificate (one candidate, one
prime, one evaluation per $b$), with which $(\star)$ is verified for all $b\equiv2,3\pmod4$ in
$[6,171]$. The uniform gap is reduced to a single $2$-adic statement.

\section{The object and the gap}
For $\lambda=(2m-b,b)$, $s=1+i$, $A=1+su+u^2$, set $g_b(m)=[u^b]A^m$,
$G_b(m)=[u^b]((1-u)A^m)=g_b-g_{b-1}$, and $I_b(m)=\Imag G_b(m)\in\Z[m]$ (degree $b$). The two-row
$d=4$ fiber law $G_{(2m-b,b)}(i)=0\iff(2,2)$ reduces, after dividing out the forced integer roots
$m=0,1,\dots,R$ ($R=\lfloor(b-1)/2\rfloor$), to the monic-over-$\Q$ polynomial $Q_b(m)$ of degree
$d=\lfloor b/2\rfloor$, and to
\[
(\star)\qquad \text{for } b\equiv2,3\!\!\pmod4,\ b\ge3:\quad Q_b(m)\text{ has no integer root }m\ge b.
\]
Classes $b\equiv0,1\pmod4$ are already proved unconditionally. Closing $(\star)$ finishes the law.

\section{The engine}
Writing $A=(su)+(1+u^2)$ and expanding gives $g_b(m)=\sum_j\binom{m}{j}s^j\binom{m-j}{(b-j)/2}$
(only $j\equiv b\bmod2$). Subtracting $g_{b-1}$:
\[
\boxed{\;I_b(m)=\sum_{j=1}^{b}(-1)^{b-j}\,\Imag(s^j)\,\binom{m}{j}\binom{m-j}{\lfloor(b-j)/2\rfloor}\;},
\quad \Imag(s^j)=\begin{cases}0,&4\mid j,\\ \{1,2,2\}[j\bmod4]\,(-4)^{\lfloor j/4\rfloor},&\text{else.}\end{cases}
\]
Verified equal to the direct extraction $\Imag[u^b]((1-u)(1+su+u^2)^m)$.

\section{The mod-2 collapse and the unique 2-adic root}
\begin{lemma}[mod-2 collapse]
$I_b(m)\equiv m\binom{m-1}{R}\pmod2$, $R=\lfloor(b-1)/2\rfloor$.
\end{lemma}
\begin{proof}
For $j\ge4$, $(-4)^{\lfloor j/4\rfloor}\equiv0\pmod2$; for $j\in\{1,2,3\}$, $\Imag(s^j)=1,2,2\equiv
1,0,0\pmod2$. So mod $2$ only $j=1$ survives:
$I_b\equiv(-1)^{b-1}\binom{m}{1}\binom{m-1}{R}\equiv m\binom{m-1}{R}$.
\end{proof}

The bare factor $m$ blocked every prior single-prime attempt (root $m\equiv0$); but it makes the
$2$-adic root \emph{unique}.

\begin{lemma}[mod-2 factorisation; prior cycle via Mahler coefficients, re-verified $6\le b\le80$]
For the monic, $2$-integral $Q_b$, \ $Q_b\equiv m\,(m^2+m+1)^{(d-1)/2}\pmod2$, with $m^2+m+1$
irreducible over $\F_2$ and $(d-1)/2=\lfloor(b-1)/4\rfloor\in\Z$.
\end{lemma}

\begin{theorem}[unique $2$-adic root]
For every $b\equiv2,3\pmod4$, $Q_b$ has exactly one root $\rho_b\in\Z_2$, and $\rho_b$ is even.
\end{theorem}
\begin{proof}
By the Lemma, over $\F_2$ the only root of $Q_b$ is $m=0$ ($m^2+m+1$ has no $\F_2$-root, its values
at $0,1$ being $1$), and it is \emph{simple}: $(Q_b\bmod2)'(0)=(m^2+m+1)^{(d-1)/2}|_{m=0}=1\ne0$.
Hensel's lemma lifts it to a unique $\rho_b\in\Z_2$ with $\rho_b\equiv0\pmod2$. Conversely any
$\Z_2$-root reduces mod $2$ to an $\F_2$-root, hence to $m=0$, hence (Hensel uniqueness) equals
$\rho_b$. So $\rho_b$ is the unique $\Z_2$-root and is even.
\end{proof}

\noindent (The $2$-adic Newton polygon $[\text{slope }-1,\text{run }1]\oplus[\text{slope }0,
\text{run }d-1]$, computed for $b\le31$, sharpens $v_2(\rho_b)=1$ and places the unit roots in the
unramified quadratic extension; only $v_2(\rho_b)\ge1$ is needed below.)

\begin{corollary}
$Q_b$ has at most one rational root, namely $\rho_b$, which is even.
\end{corollary}
\begin{proof}
A rational root is an integer (monic), hence lies in $\Z\subset\Z_2$ and is a root there; by the
theorem it equals $\rho_b$.
\end{proof}

This inverts the obstruction: the bare $m$ collapses the candidate set from $\sim0.1\,b^2$ integers
to a \emph{single} number $\rho_b$, computable from the single prime $2$.

\section{The single-candidate certificate}
\begin{theorem}[certificate]
Fix $b\equiv2,3\pmod4$. Let $U_b=1+\max_i|a_i|$ be the Cauchy bound on the roots of monic
$Q_b=\sum a_i m^i$. Choose $K$ with $2^K>U_b$ and lift the $2$-adic root to
$r:=\rho_b\bmod2^K\in[0,2^K)$. Then the only possible integer root $m_0\in[b,U_b]$ is $m_0=r$; if
$r\notin[b,U_b]$ or $Q_b(r)\ne0$, then $Q_b$ has no integer root $\ge b$, i.e.\ $(\star)$ holds for $b$.
\end{theorem}
\begin{proof}
An integer root $m_0\ge b$ has $0<m_0\le U_b<2^K$ and $m_0=\rho_b$ in $\Z_2$, so $m_0\equiv\rho_b
\pmod{2^K}$; its unique representative in $[0,2^K)$ is $r$, whence $m_0=r$. One exact evaluation
$Q_b(r)$ settles it.
\end{proof}

\begin{theorem}[frontier]
$(\star)$ holds for all $b\equiv2,3\pmod4$ with $6\le b\le171$; hence the two-row $d=4$ law holds
for these $b$. Cross-checked against exact rational-root enumeration for $b\le40$.
\end{theorem}

\noindent Examples: $b=6$, $U_6=54$, sole candidate $r=18$, $Q_6(18)=1920\ne0$; $b=7$, $U_7=791$,
sole candidate $r=610$, $Q_7(610)\ne0$.

\section{The residual gap}
\begin{quote}
\textbf{Open ($(\star)$, uniform).} For all $b\equiv2,3\pmod4$: the unique $2$-adic root
$\rho_b\in\Z_2$ is not a rational integer in $[b,\infty)$. Equivalently, $Q_b$ has no linear factor
over $\Q$, i.e.\ irreducibility (or just no-linear-factor) of the Gegenbauer-in-parameter family
$g_b(m)=C_b^{(-m)}(-(1+i)/2)$.
\end{quote}
By the corollary this is a statement about \emph{one} $2$-adic number per $b$. For irreducible
$Q_b$ (empirically all, $b\le171$) it is automatic, and Jordan's derangement theorem $+$ Chebotarev
furnish a positive density of rootless primes (finite certificates). Rigorously dead and not to be
retried: single-prime Newton polygon (integral slopes at $2$; no full non-integral run for
$p\le7$), uniform witness prime, finite covering set, analytic/log-concavity domination (in-range
real roots), and the $q$-Lucas / Gaussian-binomial shortcut.

\end{document}
